% Options for packages loaded elsewhere
\PassOptionsToPackage{unicode}{hyperref}
\PassOptionsToPackage{hyphens}{url}
\documentclass[
  ignorenonframetext,
]{beamer}
\newif\ifbibliography
\usepackage{pgfpages}
\setbeamertemplate{caption}[numbered]
\setbeamertemplate{caption label separator}{: }
\setbeamercolor{caption name}{fg=normal text.fg}
\beamertemplatenavigationsymbolsempty
% remove section numbering
\setbeamertemplate{part page}{
  \centering
  \begin{beamercolorbox}[sep=16pt,center]{part title}
    \usebeamerfont{part title}\insertpart\par
  \end{beamercolorbox}
}
\setbeamertemplate{section page}{
  \centering
  \begin{beamercolorbox}[sep=12pt,center]{section title}
    \usebeamerfont{section title}\insertsection\par
  \end{beamercolorbox}
}
\setbeamertemplate{subsection page}{
  \centering
  \begin{beamercolorbox}[sep=8pt,center]{subsection title}
    \usebeamerfont{subsection title}\insertsubsection\par
  \end{beamercolorbox}
}
% Prevent slide breaks in the middle of a paragraph
\widowpenalties 1 10000
\raggedbottom
\AtBeginPart{
  \frame{\partpage}
}
\AtBeginSection{
  \ifbibliography
  \else
    \frame{\sectionpage}
  \fi
}
\AtBeginSubsection{
  \frame{\subsectionpage}
}
\usepackage{iftex}
\ifPDFTeX
  \usepackage[T1]{fontenc}
  \usepackage[utf8]{inputenc}
  \usepackage{textcomp} % provide euro and other symbols
\else % if luatex or xetex
  \usepackage{unicode-math} % this also loads fontspec
  \defaultfontfeatures{Scale=MatchLowercase}
  \defaultfontfeatures[\rmfamily]{Ligatures=TeX,Scale=1}
\fi
\usepackage{lmodern}
\ifPDFTeX\else
  % xetex/luatex font selection
\fi
% Use upquote if available, for straight quotes in verbatim environments
\IfFileExists{upquote.sty}{\usepackage{upquote}}{}
\IfFileExists{microtype.sty}{% use microtype if available
  \usepackage[]{microtype}
  \UseMicrotypeSet[protrusion]{basicmath} % disable protrusion for tt fonts
}{}
\makeatletter
\@ifundefined{KOMAClassName}{% if non-KOMA class
  \IfFileExists{parskip.sty}{%
    \usepackage{parskip}
  }{% else
    \setlength{\parindent}{0pt}
    \setlength{\parskip}{6pt plus 2pt minus 1pt}}
}{% if KOMA class
  \KOMAoptions{parskip=half}}
\makeatother
\usepackage{color}
\usepackage{fancyvrb}
\newcommand{\VerbBar}{|}
\newcommand{\VERB}{\Verb[commandchars=\\\{\}]}
\DefineVerbatimEnvironment{Highlighting}{Verbatim}{commandchars=\\\{\}}
% Add ',fontsize=\small' for more characters per line
\newenvironment{Shaded}{}{}
\newcommand{\AlertTok}[1]{\textcolor[rgb]{1.00,0.00,0.00}{\textbf{#1}}}
\newcommand{\AnnotationTok}[1]{\textcolor[rgb]{0.38,0.63,0.69}{\textbf{\textit{#1}}}}
\newcommand{\AttributeTok}[1]{\textcolor[rgb]{0.49,0.56,0.16}{#1}}
\newcommand{\BaseNTok}[1]{\textcolor[rgb]{0.25,0.63,0.44}{#1}}
\newcommand{\BuiltInTok}[1]{\textcolor[rgb]{0.00,0.50,0.00}{#1}}
\newcommand{\CharTok}[1]{\textcolor[rgb]{0.25,0.44,0.63}{#1}}
\newcommand{\CommentTok}[1]{\textcolor[rgb]{0.38,0.63,0.69}{\textit{#1}}}
\newcommand{\CommentVarTok}[1]{\textcolor[rgb]{0.38,0.63,0.69}{\textbf{\textit{#1}}}}
\newcommand{\ConstantTok}[1]{\textcolor[rgb]{0.53,0.00,0.00}{#1}}
\newcommand{\ControlFlowTok}[1]{\textcolor[rgb]{0.00,0.44,0.13}{\textbf{#1}}}
\newcommand{\DataTypeTok}[1]{\textcolor[rgb]{0.56,0.13,0.00}{#1}}
\newcommand{\DecValTok}[1]{\textcolor[rgb]{0.25,0.63,0.44}{#1}}
\newcommand{\DocumentationTok}[1]{\textcolor[rgb]{0.73,0.13,0.13}{\textit{#1}}}
\newcommand{\ErrorTok}[1]{\textcolor[rgb]{1.00,0.00,0.00}{\textbf{#1}}}
\newcommand{\ExtensionTok}[1]{#1}
\newcommand{\FloatTok}[1]{\textcolor[rgb]{0.25,0.63,0.44}{#1}}
\newcommand{\FunctionTok}[1]{\textcolor[rgb]{0.02,0.16,0.49}{#1}}
\newcommand{\ImportTok}[1]{\textcolor[rgb]{0.00,0.50,0.00}{\textbf{#1}}}
\newcommand{\InformationTok}[1]{\textcolor[rgb]{0.38,0.63,0.69}{\textbf{\textit{#1}}}}
\newcommand{\KeywordTok}[1]{\textcolor[rgb]{0.00,0.44,0.13}{\textbf{#1}}}
\newcommand{\NormalTok}[1]{#1}
\newcommand{\OperatorTok}[1]{\textcolor[rgb]{0.40,0.40,0.40}{#1}}
\newcommand{\OtherTok}[1]{\textcolor[rgb]{0.00,0.44,0.13}{#1}}
\newcommand{\PreprocessorTok}[1]{\textcolor[rgb]{0.74,0.48,0.00}{#1}}
\newcommand{\RegionMarkerTok}[1]{#1}
\newcommand{\SpecialCharTok}[1]{\textcolor[rgb]{0.25,0.44,0.63}{#1}}
\newcommand{\SpecialStringTok}[1]{\textcolor[rgb]{0.73,0.40,0.53}{#1}}
\newcommand{\StringTok}[1]{\textcolor[rgb]{0.25,0.44,0.63}{#1}}
\newcommand{\VariableTok}[1]{\textcolor[rgb]{0.10,0.09,0.49}{#1}}
\newcommand{\VerbatimStringTok}[1]{\textcolor[rgb]{0.25,0.44,0.63}{#1}}
\newcommand{\WarningTok}[1]{\textcolor[rgb]{0.38,0.63,0.69}{\textbf{\textit{#1}}}}
\usepackage{longtable,booktabs,array}
\newcounter{none} % for unnumbered tables
\usepackage{calc} % for calculating minipage widths
\usepackage{caption}
% Make caption package work with longtable
\makeatletter
\def\fnum@table{\tablename~\thetable}
\makeatother
\setlength{\emergencystretch}{3em} % prevent overfull lines
\providecommand{\tightlist}{%
  \setlength{\itemsep}{0pt}\setlength{\parskip}{0pt}}
% Slide layout defaults for warning-clean scientific decks.
\usepackage{etoolbox}
\IfFileExists{xurl.sty}{\usepackage{xurl}}{}
\IfFileExists{seqsplit.sty}{\usepackage{seqsplit}}{\newcommand{\seqsplit}[1]{#1}}
\protected\def\breakseq#1{\seqsplit{#1}}
\protected\def\breaktt#1{\begingroup\ttfamily\seqsplit{#1}\endgroup}
\setlength{\emergencystretch}{6em}
\tolerance=5000
\hbadness=10000
\hfuzz=1pt
\setlength{\tabcolsep}{2pt}
\AtBeginEnvironment{longtable}{\tiny\renewcommand{\arraystretch}{0.86}\setlength{\tabcolsep}{1pt}}
\AtBeginEnvironment{tabular}{\tiny\renewcommand{\arraystretch}{0.86}\setlength{\tabcolsep}{1pt}}

% Natbib and cross-reference fallbacks — slides don't load natbib
% or cleveref, but combined-PDF manuscript prose may emit these
% commands. The fallback renders citations as a bracketed key list
% and cross-references as detokenized labels so slides don't fail on
% undefined control sequences or raw underscores. \providecommand is
% a no-op if packages load later.
\providecommand{\citep}[1]{[#1]}
\providecommand{\citet}[1]{#1}
\providecommand{\citealp}[1]{#1}
\providecommand{\citeauthor}[1]{#1}
\providecommand{\citeyear}[1]{#1}
\providecommand{\cref}[1]{\texttt{\detokenize{#1}}}
\providecommand{\Cref}[1]{\texttt{\detokenize{#1}}}

\newtheorem{proposition}{Proposition}
\newtheorem{hypothesis}{Hypothesis}
\usepackage{bookmark}
\IfFileExists{xurl.sty}{\usepackage{xurl}}{} % add URL line breaks if available
\urlstyle{same}
% fallback for those not using the hyperref driver hyperxmp:
\makeatletter
\@ifundefined{xmpquote}{\newcommand{\xmpquote}[1]{#1}}{}
\makeatother
\hypersetup{
  hidelinks,
  pdfcreator={LaTeX via pandoc}}

\author{\texorpdfstring{}{}}
\date{}

\begin{document}

\section{Reproducibility}\label{sec:reproducibility}

\begin{frame}[fragile,allowframebreaks]{Reproducibility}
Reproducibility in computational research has well-documented
prerequisites: open data, open code, and a deterministic build that can
be re-run from scratch \breakseq{[@peng2011reproducible]}. The bundled
\breaktt{manuscript/config.yaml} is intentionally configured to satisfy
all three for \textbf{strict reproducibility}:

\begin{enumerate}
\tightlist
\item
  \texttt{search.sources:\ {[}local{]}} consumes
  \breaktt{data/corpus.json}, which is a curated and committed JSON
  corpus. No network is required to run the pipeline.
\item
  \breaktt{search.cache\_dir:\ output/search/cache} writes deterministic
  JSON cache files; running the same query twice produces a
  byte-identical artifact tree (modulo timestamp metadata in the cache
  file itself).
\item
  \breaktt{enrichment.fetch\_abstracts:\ true} reads abstracts directly
  from the corpus when present; no network fetch is required.
\item
  \breaktt{enrichment.fetch\_fulltext:\ false} is the default ---
  full-text fetching is opt-in and gated behind the optional
  \texttt{pypdf} dependency (\breaktt{uv\ sync\ -\/-group\ rendering}).
\item
  \breaktt{llm.enabled:\ false} is the default --- the LLM stage is
  opt-in and requires a running \texttt{ollama\ serve}. When enabled,
  \texttt{seed:\ 42} and \breaktt{temperature:\ 0.0} are pinned.
\end{enumerate}
\end{frame}

\begin{frame}[fragile,allowframebreaks]{Switching to live search}
\protect\phantomsection\label{switching-to-live-search}
Replace the \texttt{sources} list with the desired backend set:

\begin{Shaded}
\begin{Highlighting}[]
\FunctionTok{search}\KeywordTok{:}
\AttributeTok{  }\FunctionTok{query}\KeywordTok{:}\AttributeTok{ }\StringTok{"your topic"}
\AttributeTok{  }\FunctionTok{sources}\KeywordTok{:}\AttributeTok{ }\KeywordTok{[}\AttributeTok{arxiv}\KeywordTok{,}\AttributeTok{ crossref}\KeywordTok{]}
\AttributeTok{  }\FunctionTok{crossref\_mailto}\KeywordTok{:}\AttributeTok{ }\StringTok{"you@example.org"}
\end{Highlighting}
\end{Shaded}

A first live run populates \breaktt{output/search/cache/}. Commit that
directory to the repo and the pipeline becomes reproducible across
machines without further configuration changes.
\end{frame}

\begin{frame}[fragile,allowframebreaks]{Determinism guarantees}
\protect\phantomsection\label{determinism-guarantees}
The full determinism contract is itemised in \breakseq{[@tbl:determinism]}:
every pipeline stage is annotated as fully, conditionally, or
non-deterministic, with an explicit mechanism column so reviewers can
audit each row independently.

\begin{longtable}[]{@{}
  >{\raggedright\arraybackslash}p{(\linewidth - 4\tabcolsep) * \real{0.3333}}
  >{\raggedright\arraybackslash}p{(\linewidth - 4\tabcolsep) * \real{0.3333}}
  >{\raggedright\arraybackslash}p{(\linewidth - 4\tabcolsep) * \real{0.3333}}@{}}
\caption{Determinism contract by pipeline stage. Cached stages are
byte-stable across reruns; live stages depend on the upstream source and
are pinned to the cache file once a successful run
completes.}\label{tbl:determinism}\tabularnewline
\toprule\noalign{}
\begin{minipage}[b]{\linewidth}\raggedright
Stage
\end{minipage} & \begin{minipage}[b]{\linewidth}\raggedright
Deterministic?
\end{minipage} & \begin{minipage}[b]{\linewidth}\raggedright
Mechanism
\end{minipage} \\
\midrule\noalign{}
\endfirsthead
\toprule\noalign{}
\begin{minipage}[b]{\linewidth}\raggedright
Stage
\end{minipage} & \begin{minipage}[b]{\linewidth}\raggedright
Deterministic?
\end{minipage} & \begin{minipage}[b]{\linewidth}\raggedright
Mechanism
\end{minipage} \\
\midrule\noalign{}
\endhead
Search (cached hit) & yes & \texttt{SearchCache} JSON files \\
Search (cache miss) & no & live API \\
Dedup / merge & yes & DOI / arXiv-id canonical keys; tie-break by score
then year \\
Citation-key generation & yes & unicode folding + stop-word skip;
collision suffix is deterministic \\
BibTeX writer & yes & byte-stable format pinning (verified by
\breaktt{tests/infra\_tests/reference/}) \\
Abstract fetch & yes (cached) / no (live) & per-paper
\texttt{\textless{}safe\_id\textgreater{}.txt} cache \\
Fulltext fetch & yes (cached) / mostly (live) & per-paper
\texttt{\textless{}safe\_id\textgreater{}.\{pdf,txt\}} cache; live
fetch's \texttt{pypdf} text extraction is not bit-stable across
versions \\
LLM synthesis & mostly & \texttt{seed=42}, \texttt{temperature=0.0};
Ollama deterministic up to its own minor variance \\
Figure generation & yes (within Matplotlib version) & fixed palette,
fixed bin width, no random subsampling \\
\bottomrule\noalign{}
\end{longtable}
\end{frame}

\begin{frame}[fragile,allowframebreaks]{Verifying reproducibility
locally}
\protect\phantomsection\label{verifying-reproducibility-locally}
\begin{Shaded}
\begin{Highlighting}[]
\CommentTok{\# Run twice; nothing in output/ should diff except the cache timestamps.}
\ExtensionTok{uv}\NormalTok{ run python projects/templates/template\_search\_project/scripts/run\_search\_pipeline.py}
\FunctionTok{mv}\NormalTok{ projects/templates/template\_search\_project/output projects/templates/template\_search\_project/output\_first}
\ExtensionTok{uv}\NormalTok{ run python projects/templates/template\_search\_project/scripts/run\_search\_pipeline.py}
\FunctionTok{diff} \AttributeTok{{-}ru} \DataTypeTok{\textbackslash{}}
\NormalTok{    projects/templates/template\_search\_project/output\_first/corpus.json }\DataTypeTok{\textbackslash{}}
\NormalTok{    projects/templates/template\_search\_project/output/corpus.json}
\end{Highlighting}
\end{Shaded}

The only expected differences are inside
\breaktt{output/search/cache/search\_*.json}, where \texttt{\_cached\_at}
is wall-clock time at write.
\end{frame}

\begin{frame}[fragile,allowframebreaks]{Limitations}
\protect\phantomsection\label{limitations}
The reproducibility contract enumerated in \breakseq{[@sec:reproducibility]}
(items 1--5 and \breakseq{[@tbl:determinism]}) does not eliminate the following
well-defined sources of non-reproducibility, which are surfaced here so
reviewers can audit them explicitly rather than inferring from the
contract table:

\begin{itemize}
\tightlist
\item
  \textbf{Live search drift.} When \breaktt{config.search.sources}
  includes \texttt{arxiv} or \texttt{crossref}, the first cache-miss
  invocation hits the live API; the cached JSON freezes that response,
  but two cold-start clones running on different days will see different
  paper sets. Pin a \texttt{LocalBackend} corpus or commit
  \breaktt{output/search/cache/} to break this dependency.
\item
  \textbf{\texttt{pypdf} version drift.} The fulltext fetcher uses
  \texttt{pypdf} to extract text from a downloaded PDF. \texttt{pypdf}'s
  text-extraction algorithm is not bit-stable across major versions;
  upgrading \texttt{pypdf} can produce different
  \texttt{\textless{}safe\_id\textgreater{}.txt} cache contents from the
  same source PDF. The PDF bytes themselves are bit-stable so the cache
  freezes the inputs, not the extraction.
\item
  \textbf{Ollama version drift.} Pinning \texttt{seed=42} and
  \texttt{temperature=0.0} controls Ollama's sampling, but the model
  weights, tokenizer, and template can change between Ollama releases.
  Document the Ollama version alongside \breaktt{config.llm.model} when
  archiving a run for replication.
\item
  \textbf{Paperclip backend status.} The \texttt{paperclip} backend is
  opt-in and currently degrades to HTTP 405 on the production endpoint;
  the run records the error in
  \texttt{SearchResult.errors{[}paperclip{]}} and continues. Treat
  \texttt{paperclip} results as advisory until the upstream service
  stabilises.
\item
  \textbf{External backend behaviour outside this project's control.}
  arXiv and Crossref are the source of truth; this project is faithful
  to whatever they return. A retraction, metadata fix, or DOI assignment
  upstream will alter the cache on the next cold-start invocation.
\end{itemize}

These limitations bound \emph{what} the cache + seed + corpus pinning
achieves. Inside those bounds, the contract in \breakseq{[@tbl:determinism]} is
total: every cached pipeline stage is byte-stable across reruns
(verified by
\breaktt{tests/test\_pipeline.py::TestRunLiteraturePipeline::test\_bibtex\_byte\_identical\_across\_reruns}).
\end{frame}

\end{document}
